\section*{Capítulo \ref{ch:g9:fractions} --- Fracciones y potencias}

\begin{solution}{exo:g9:fractions:1}
$\dfrac13 + \dfrac15 = \dfrac{5}{15} + \dfrac{3}{15} = \dfrac{8}{15}$.

$\dfrac72 - \dfrac53 = \dfrac{21}{6} - \dfrac{10}{6} = \dfrac{11}{6}$.

$\dfrac{5}{12} + \dfrac38 = \dfrac{10}{24} + \dfrac{9}{24} =
\dfrac{19}{24}$.

$3 - \dfrac45 = \dfrac{15}{5} - \dfrac45 = \dfrac{11}{5}$.
\end{solution}

\begin{solution}{exo:g9:fractions:2}
$\dfrac59 \times \dfrac{27}{35} = \dfrac{5 \times 27}{9 \times 35}
= \dfrac{27}{9} \times \dfrac{5}{35} = 3 \times \dfrac17 = \dfrac37$
(simplificando por $9$ y por $5$).

$\dfrac{8}{15} \div \dfrac45 = \dfrac{8}{15} \times \dfrac54
= \dfrac{8 \times 5}{15 \times 4} = \dfrac{40}{60} = \dfrac23$.

$\dfrac23 \times \dfrac34 \times \dfrac45 = \dfrac{2 \times 3 \times 4}
{3 \times 4 \times 5} = \dfrac25$ (los $3$ y los $4$ se cancelan).
\end{solution}

\begin{solution}{exo:g9:fractions:3}
$2^5 = 32$; \quad $(-2)^4 = 16$; \quad $-2^4 = -16$; \quad
$4^{-2} = \frac{1}{16}$; \quad
$\left(\frac23\right)^3 = \frac{8}{27}$; \quad $7^0 = 1$.
\end{solution}

\begin{solution}{exo:g9:fractions:4}
$5^3 \times 5^6 = 5^9$; \quad $\dfrac{7^9}{7^5} = 7^4$; \quad
$\left(2^3\right)^4 = 2^{12}$; \quad
$\dfrac{3^2 \times 3^7}{3^5} = 3^{9-5} = 3^4$; \quad
$2^5 \times 4 = 2^5 \times 2^2 = 2^7$.
\end{solution}

\begin{solution}{exo:g9:fractions:5}
$45\,000 = 4.5 \times 10^4$; \quad
$0.0072 = 7.2 \times 10^{-3}$; \quad
$302.5 = 3.025 \times 10^2$; \quad
$0.000\,000\,9 = 9 \times 10^{-7}$; \quad
doce millones $= 1.2 \times 10^7$.
\end{solution}

\begin{solution}{exo:g9:fractions:6}
$(2 \times 10^7) \times (4.5 \times 10^{-3}) = 9 \times 10^{4}$.

$\dfrac{9 \times 10^5}{3 \times 10^{-2}} = 3 \times 10^{5 - (-2)}
= 3 \times 10^7$.

$(5 \times 10^3)^2 = 25 \times 10^6 = 2.5 \times 10^7$.
\end{solution}

\begin{solution}{exo:g9:fractions:7}
El agua añadida representa
$\dfrac34 - \dfrac35 = \dfrac{15}{20} - \dfrac{12}{20} = \dfrac{3}{20}$
del depósito. Así que $\frac{3}{20}$ de la \omterm{def:g2:measure:units}{capacidad} son $30$ litros:
$\frac{1}{20}$ son $10$ litros, y la \omterm{def:g2:measure:units}{capacidad} es de $200$ litros.
\end{solution}

\begin{solution}{exo:g9:fractions:8}
Después de Alicia queda $1 - \frac14 = \frac34$ del pastel. Ben se come
$\frac13 \times \frac34 = \frac14$, y quedan
$\frac34 - \frac14 = \frac12$. Carla se come la \omterm{def:g3:division:half}{mitad} de eso, $\frac14$,
y queda $\frac14$ del pastel.
\end{solution}

\begin{solution}{exo:g9:fractions:9}
Tiempo $= \dfrac{\text{distancia}}{\text{velocidad}} =
\dfrac{1.5 \times 10^{11}}{3 \times 10^8} = 0.5 \times 10^3 = 500$
segundos, es decir $8$ minutos y $20$ segundos.
\end{solution}

\begin{solution}{exo:g9:fractions:10}
Escribe los dos números como \omterm{def:g9:fractions:power}{potencias} de \omterm{def:g8:powers:def}{exponente} $10$:
\[
2^{30} = \left(2^3\right)^{10} = 8^{10},
\qquad
3^{20} = \left(3^2\right)^{10} = 9^{10}.
\]
Como $8 < 9$, al multiplicar diez copias de cada uno: $8^{10} < 9^{10}$, luego
$2^{30} < 3^{20}$.
\end{solution}

\begin{solution}{pb:g9:fractions:1}
\textbf{1.} $\frac14 = 0.25$ (se para); $\frac13 = 0.333\ldots$
(periodo $3$); $\frac56 = 0.8333\ldots$ (periodo $3$ tras el
$8$); $\frac{3}{11} = 0.272727\ldots$ (periodo $27$).

\textbf{2.} En cada paso de la división entre $b$, el \omterm{def:g3:division:remainder}{resto}
es uno de los $b$ valores $0, 1, \dots, b - 1$
(\cref{thm:g6:wholes:division}). Si aparece un \omterm{def:g3:division:remainder}{resto} $0$, la
división se para. Si no, tras $b$ pasos como mucho algún
\omterm{def:g3:division:remainder}{resto} tiene que aparecer por segunda vez --- hay más pasos
que \omterm{def:g3:division:remainder}{restos} posibles --- y a partir de ese momento la división
repite exactamente la sucesión de cifras que produjo tras la
primera aparición: la escritura decimal cicla para siempre. Pararse o
repetirse: no existe un tercer comportamiento.

\textbf{3.} $\frac17 = 0.142857\,142857\ldots$: el periodo
$142857$ tiene longitud $6$. La pregunta 2 prometía una repetición en
$7$ pasos como mucho --- aquí los \omterm{def:g3:division:remainder}{restos}
$1, 3, 2, 6, 4, 5$ aparecen todos antes de que vuelva el $1$: la
cantinela más larga posible para una división entre $7$.

\textbf{4.} $\frac ab$ se para exactamente cuando algún
$\frac{a \times k}{b \times k}$ tiene \omterm{def:g4:fractions:def}{denominador} $10$, $100$,
$1000, \dots$ --- es decir, cuando el \omterm{def:g4:fractions:def}{denominador} se puede
multiplicar hasta una \omterm{def:g9:fractions:power}{potencia} de diez, como $4 \times 25 = 100$ o
$8 \times 125 = 1\,000$. Para $3$ y $6$ no hay nada que hacer: una
\omterm{def:g9:fractions:power}{potencia} de diez tiene \omterm{def:g1:addition:def}{suma} de cifras $1$, y todo \omterm{def:g5:division:multiple}{múltiplo} de $3$ tiene
una \omterm{def:g1:addition:def}{suma} de cifras que es \omterm{def:g5:division:multiple}{múltiplo} de $3$ --- así que ningún \omterm{def:g5:division:multiple}{múltiplo} de $3$
(y por tanto ninguno de $6$) es nunca $10$, $100$, $1\,000, \dots$ Para
$11$: dividir $10$, $100$, $1\,000$, $10\,000$ entre $11$ siempre
deja \omterm{def:g3:division:remainder}{resto} ($10$, $1$, $10$, $1, \dots$), nunca $0$.
Los \omterm{def:g4:fractions:def}{denominadores} que alcanzan una \omterm{def:g9:fractions:power}{potencia} de diez se paran; todos los demás se repiten.

\textbf{5.} $\frac{9}{40}$: $40 \times 25 = 1\,000$, se para:
$\frac{225}{1000} = 0.225$. $\frac{33}{50}$:
$50 \times 2 = 100$, se para: $\frac{66}{100} = 0.66$.
$\frac{5}{12}$: $12$ es \omterm{def:g5:division:multiple}{múltiplo} de $3$, se repite
($0.41666\ldots$).

\textbf{6.} $10x = 7.777\ldots$, luego $10x - x = 7$ (las colas
se cancelan a la perfección), $9x = 7$ y $x = \frac79$. Comprobación por división:
$7 \div 9 = 0.777\ldots$

\textbf{7.} El periodo tiene dos cifras, así que desplaza con $100$:
$100x - x = 72.7272\ldots - 0.7272\ldots = 72$, de donde
$99x = 72$ y $x = \frac{72}{99} = \frac{8}{11}$.

\textbf{8.} $1000x = 583.333\ldots$ y $100x = 58.333\ldots$;
restando, $900x = 525$, luego
$x = \frac{525}{900} = \frac{7}{12}$ (divide entre $75$). En efecto,
$\frac{7}{12} = 0.58333\ldots$

\textbf{9.} $10x - x = 9.999\ldots - 0.999\ldots = 9$, luego
$9x = 9$ y $x = 1$. No está «justo por debajo» de $1$:
$0.999\ldots$ \emph{es} el número $1$, escrito con un segundo
disfraz. Cualquier número estrictamente menor que $1$ deja un hueco, y
\cref{pb:g6:decimals:1} mostró que un hueco siempre contiene más
números --- mientras que entre $0.999\ldots$ y $1$ no cabe nada.
Los nueves sin fin son el \omterm{def:g3:division:remainder}{resto} de chocolate del
\cref{pb:g6:fractions:1} encogido por debajo de cualquier cantidad positiva:
no queda nada.

\textbf{10.} $2.4545\ldots = 2 + \frac{45}{99} = 2 +
\frac{5}{11} = \frac{27}{11}$. Y
$0.123123\ldots = \frac{123}{999} = \frac{41}{333}$ (divide entre
$3$).

\textbf{11.} Diccionario: $0.037037\ldots = \frac{37}{999}$.
Como $999 = 27 \times 37$, esto se simplifica a $\frac{1}{27}$:
la \omterm{def:g9:fractions:fraction}{fracción} $\frac{1}{27}$ canturrea «$037$». Comprobación por división:
$1 \div 27 = 0.037037\ldots$

\textbf{12.} $0.0272727\ldots = \frac{1}{10} \times
0.272727\ldots = \frac{1}{10} \times \frac{27}{99}
= \frac{27}{990} = \frac{3}{110}$ (divide entre $9$).

\textbf{13.} Los bloques de \omterm{def:g1:counting:zero}{ceros} crecen sin fin, así que ningún bloque
fijo puede repetirse para siempre: \emph{no} es un decimal periódico
(y tampoco se para nunca). Por la Parte I, toda \omterm{def:g9:fractions:fraction}{fracción} se para o
se repite --- así que este número no es ninguna \omterm{def:g9:fractions:fraction}{fracción}: es
irracional, hecho a medida. La superestrella de la irracionalidad,
$\sqrt2$, queda desenmascarada en el \cref{ch:g9:sqrt}.

\textbf{14.} $2^{30} = \left(2^{10}\right)^3 \approx
(1.024)^3 \times 10^9 \approx 1.07 \times 10^9$: un número
un poco mayor que $10^9$, luego con $10$ cifras. (Exactamente:
$2^{30} = 1\,073\,741\,824$.)

\textbf{15.} $0.20262026\ldots = \frac{2026}{9999}$, ya
irreducible ($9999 = 3 \times 3 \times 11 \times 101$ no comparte
ningún factor con $2026 = 2 \times 1013$). El diccionario
completo: \emph{las fracciones son exactamente los decimales que
se paran o se repiten} --- los \omterm{def:g4:fractions:def}{denominadores} amigos del diez se paran, todos
los demás canturrean un periodo; y, recíprocamente, toda cantinela, incluso tu
año de graduación, es una \omterm{def:g9:fractions:fraction}{fracción} disfrazada.
\end{solution}
